The bug report is valid: the statement “no integer multiplied by \(2\) equals \(1\)” should be justified. Here is a revised proof that explicitly proves that fact before using it.

```latex
\begin{proof}
Recall that an integer \(m\) is even if there exists an integer \(t\) such that
\(m=2t\).

We first record a basic fact: there is no integer \(k\) such that
\[
2k=1.
\]
Indeed, let \(k\in\mathbb{Z}\). Either \(k\leq 0\) or \(k>0\). If \(k\leq 0\),
then \(2k\leq 0\), so \(2k\neq 1\). If \(k>0\), then, since \(k\) is an
integer, \(k\geq 1\), and hence \(2k\geq 2\), so again \(2k\neq 1\). Therefore
no integer \(k\) satisfies \(2k=1\).

We now prove the parity fact that will be used twice: if \(n^2\) is even, then
\(n\) is even.

Let \(n\in\mathbb{Z}\). By the division algorithm with divisor \(2\), there
exist integers \(q\) and \(r\), with \(r\in\{0,1\}\), such that
\[
n=2q+r.
\]
Suppose \(n\) is not even. If \(r=0\), then \(n=2q\), which would make \(n\)
even. Hence \(r\neq 0\). Since \(r\in\{0,1\}\), we must have \(r=1\), so
\[
n=2q+1.
\]
Then
\[
n^2=(2q+1)^2=4q^2+4q+1=2(2q^2+2q)+1.
\]
If \(n^2\) were even, then there would exist an integer \(u\) such that
\[
n^2=2u.
\]
Combining this with the previous expression for \(n^2\), we would get
\[
2u=2(2q^2+2q)+1.
\]
Thus
\[
1=2\bigl(u-(2q^2+2q)\bigr).
\]
Since \(u-(2q^2+2q)\in\mathbb{Z}\), this contradicts the fact proved above that
no integer \(k\) satisfies \(2k=1\). Therefore, if \(n\) is not even, then
\(n^2\) is not even. Equivalently, if \(n^2\) is even, then \(n\) is even.

Now suppose, for the sake of contradiction, that \(\sqrt{2}\) is rational. Then
there exist integers \(p\) and \(q\), with \(q\neq 0\), such that
\[
\sqrt{2}=\frac{p}{q}.
\]
Consider the set
\[
S=\left\{ d\in \mathbb{Z}_{>0} : \text{there exists } c\in\mathbb{Z}
\text{ such that } \sqrt{2}=\frac{c}{d} \right\}.
\]
This set is nonempty. Indeed, if \(q>0\), then \(q\in S\). If \(q<0\), then
\(-q>0\), and
\[
\frac{p}{q}=\frac{-p}{-q},
\]
so \(-q\in S\).

By the well-ordering principle, every nonempty subset of \(\mathbb{Z}_{>0}\)
has a least element. Therefore \(S\) has a least element. Let \(b\) be the least
element of \(S\). Since \(b\in S\), there exists an integer \(a\) such that
\[
\sqrt{2}=\frac{a}{b}.
\]
By construction, \(b>0\), and \(b\) is minimal among all positive integer
denominators appearing in representations of \(\sqrt{2}\) as a ratio of
integers.

Squaring both sides gives
\[
2=\frac{a^2}{b^2}.
\]
Since \(b>0\), we have \(b\neq 0\), so multiplying both sides by \(b^2\) yields
\[
a^2=2b^2.
\]
Thus \(a^2\) is even. By the parity fact proved above, \(a\) is even, so there
exists an integer \(c\) such that
\[
a=2c.
\]
Substituting this into \(a^2=2b^2\), we get
\[
(2c)^2=2b^2.
\]
Hence
\[
4c^2=2b^2.
\]
Dividing both sides by \(2\), we obtain
\[
2c^2=b^2.
\]
Therefore \(b^2\) is even. By the same parity fact, \(b\) is even, so there
exists an integer \(d\) such that
\[
b=2d.
\]
Since \(b>0\), we have \(2d>0\), and therefore \(d>0\). Also,
\[
\sqrt{2}=\frac{a}{b}=\frac{2c}{2d}=\frac{c}{d}.
\]
Thus \(d\in S\). But \(b=2d\) and \(d>0\), so
\[
b-d=2d-d=d>0,
\]
which means
\[
d<b.
\]
This contradicts the minimality of \(b\) as the least element of \(S\).

Therefore the assumption that \(\sqrt{2}\) is rational is false. Hence
\(\sqrt{2}\) is irrational.
\end{proof}
```